\chapter{Scalar-Fiber Canonical and Resolved Closure}
\label{ch:canonical}

\section{ADM variables}

Write the spacetime metric as
\begin{equation}
\dd s^2=-N^2\dd t^2+q_{ij}(\dd x^i+N^i\dd t)(\dd x^j+N^j\dd t).
\end{equation}
Let $\pi_\phi$ and $\Pi$ be the momenta conjugate to $\phi$ and $\chi$.  The scalar-fiber matter constraints are
\begin{equation}
\boxed{
\begin{aligned}
\Hh_\perp^{\mathrm{sf}}={}&
\frac{1}{2\sqrt q}\left(\pi_\phi^2+\ip{\Pi}{\Pi}_H\right)\\
&+\frac{\sqrt q}{2}\left[
q^{ij}\partial_i\phi\partial_j\phi
+q^{ij}\ip{\partial_i\chi}{\partial_j\chi}_H
+2V_H(\phi,\chi)
\right],
\end{aligned}}
\label{eq:adm-hmatter}
\end{equation}
\begin{equation}
\boxed{
\Hh_i^{\mathrm{sf}}
=\pi_\phi\partial_i\phi
+\ip{\Pi}{\partial_i\chi}_H.}
\label{eq:adm-dmatter}
\end{equation}
The smeared generators are
\begin{equation}
H_{\mathrm{sf}}[N]=\int_\Sigma N\Hh_\perp^{\mathrm{sf}}\dd^3x,
\qquad
D_{\mathrm{sf}}[\bm N]=\int_\Sigma N^i\Hh_i^{\mathrm{sf}}\dd^3x.
\end{equation}

\section{Matter brackets and their metric term}

\begin{theorem}[Matter momentum and Hamiltonian brackets]
On a smooth spatial chart with boundary conditions that remove surface terms, the constraints \eqref{eq:adm-hmatter}--\eqref{eq:adm-dmatter} satisfy
\begin{align}
\PB{D_{\mathrm{sf}}[\bm N]}{D_{\mathrm{sf}}[\bm M]}
&=D_{\mathrm{sf}}[[\bm N,\bm M]],
\label{eq:adm-dd}\\
\PB{H_{\mathrm{sf}}[N]}{H_{\mathrm{sf}}[M]}
&=D_{\mathrm{sf}}[\bm S(N,M)],
\label{eq:adm-hh}
\end{align}
where
\begin{equation}
S^i(N,M)=q^{ij}(N\partial_jM-M\partial_jN).
\label{eq:adm-shift}
\end{equation}
The result holds for every differentiable scalar potential $V_H$, including the scale-covariant fiber interaction.  The pure-matter mixed bracket is instead
\begin{equation}
\boxed{
\begin{aligned}
\PB{D_{\mathrm{sf}}[\bm N]}{H_{\mathrm{sf}}[M]}
={}&H_{\mathrm{sf}}[\Lie_{\bm N}M]\\
&+\int_\Sigma
\frac{\delta H_{\mathrm{sf}}[M]}{\delta q_{ij}}
\Lie_{\bm N}q_{ij}\,\dd^3x.
\end{aligned}}
\label{eq:adm-dh-metric}
\end{equation}
\end{theorem}

\begin{proof}
The momentum generator acts on $\phi$ and $\chi$ by their scalar Lie derivatives and on the momenta by the corresponding density Lie derivatives.  This gives \eqref{eq:adm-dd}.  It does not act on the metric appearing in $H_{\mathrm{sf}}$.  Restoring that missing metric transformation gives the second term in \eqref{eq:adm-dh-metric}.

For the Hamiltonian bracket, vary the kinetic and spatial-gradient terms.  Terms containing $\partial V_H/\partial\phi$ and the fiber gradient $\partial V_H/\partial\chi$ appear multiplied by $NM$ in both bracket orders and cancel by antisymmetry.  The remaining cross terms are
\begin{equation}
\int_\Sigma
(N\partial_iM-M\partial_iN)q^{ij}
\left(
\pi_\phi\partial_j\phi+\ip{\Pi}{\partial_j\chi}_H
\right)\dd^3x,
\end{equation}
which is \eqref{eq:adm-hh}.
\end{proof}

The potential independence of \eqref{eq:adm-hh} is important.  Canonical closure tests covariance of the matter action.  It does not determine the numerical value of $\gamma$ or validate the production state.

\section{Closure after coupling to gravity}

Let $H_{\mathrm G},D_{\mathrm G}$ denote the gravitational constraints on a regular ADM chart.  The total generators are
\begin{equation}
H_{\mathrm{tot}}=H_{\mathrm G}+H_{\mathrm{sf}},
\qquad
D_{\mathrm{tot}}=D_{\mathrm G}+D_{\mathrm{sf}}.
\end{equation}
\begin{theorem}[Total hypersurface deformation]
On every regular ADM chart where the gravitational bracket is defined, minimal coupling gives
\begin{align}
\PB{D_{\mathrm{tot}}[\bm N]}{D_{\mathrm{tot}}[\bm M]}
&=D_{\mathrm{tot}}[[\bm N,\bm M]],
\label{eq:adm-total-dd}\\
\PB{D_{\mathrm{tot}}[\bm N]}{H_{\mathrm{tot}}[M]}
&=H_{\mathrm{tot}}[\Lie_{\bm N}M],
\label{eq:adm-total-dh}\\
\PB{H_{\mathrm{tot}}[N]}{H_{\mathrm{tot}}[M]}
&=D_{\mathrm{tot}}[\bm S(N,M)].
\label{eq:adm-total-hh}
\end{align}
\end{theorem}

\begin{proof}
Minimal coupling makes $\Hh_\perp^{\mathrm{sf}}$ a scalar density of weight one under the total momentum generator.  The metric-variation term in \eqref{eq:adm-dh-metric} is supplied by the gravitational action on $q_{ij}$.  The mixed metric variations in the total Hamiltonian bracket combine with the two pure brackets, yielding the displayed algebra.
\end{proof}

\begin{statusbox}[Closure boundary]
This result is a classical constraint theorem.  It is not a quantum anomaly theorem.  It also does not turn a nonlocal spectral derivative into a strictly finite-range microscopic lattice action.
\end{statusbox}

\section{Alias-free resolved realization}

On the periodic line, let $\mathcal T_K$ denote trigonometric polynomials of band at most $K$.  Sampling on $N_x$ points is injective below Nyquist.  If every product generated in one bracket remains below $N_x/2$, the sampled derivative, product, and normalized trace coincide with their continuum values.  The proof of the matter $H$--$H$ identity \eqref{eq:adm-hh} then transfers term by term.  The resolved test below does not test the mixed bracket or the dynamical metric sector.

For the one-dimensional scalar-fiber certificate, the Hamiltonian density is
\begin{equation}
\Hh^{\mathrm{sf}}
=\frac12\left(\pi_\phi^2+\ip{\Pi}{\Pi}_H\right)
+\frac12\left((\delta\phi)^2+\ip{\delta\chi}{\delta\chi}_H\right)
+V_H.
\label{eq:adm-1d-density}
\end{equation}
The exact target identity is
\begin{equation}
\PB{H_N[n]}{H_N[m]}
=D_N[n\delta m-m\delta n].
\label{eq:adm-resolved-identity}
\end{equation}

\begin{proposition}[Resolved matter Hamiltonian bracket]
If the fields and smearings lie in a band for which every product in the variation and bracket is resolved without aliasing, equation \eqref{eq:adm-resolved-identity} holds exactly in the sampled algebra.
\end{proposition}

\begin{proof}
The alias-free product lemma identifies every sampled multiplication and derivative in the bracket with the corresponding trigonometric-polynomial operation.  The continuum proof of \eqref{eq:adm-hh} therefore applies without a truncation remainder.
\end{proof}

The deterministic test used $N_x=128$ and band $K=3$.  Across twenty random trials, the maximum relative residual was
\begin{equation}
\boxed{\epsilon_{\mathrm{resolved}}=\BridgeSpectralResidual.}
\label{eq:adm-resolved-residual}
\end{equation}

\section{Negative controls}

Three changes break the hypotheses in distinguishable ways.
\begin{enumerate}[leftmargin=2.2em]
\item Raising the band to $K=18$ on $N_x=64$ aliases products.  The residual becomes $\BridgeAliasResidual$.
\item Replacing the spectral derivative by a centered difference destroys the exact product rule.  The residual becomes $\BridgeFiniteDifferenceResidual$.
\item Removing the detail spatial-gradient term while retaining its momentum contribution makes the Hamiltonian and momentum generators inconsistent.  The residual becomes $\BridgeMissingGradientResidual$.
\end{enumerate}

\begin{figure}[htbp]
\centering
\includegraphics[width=0.93\linewidth]{constraint_negative_controls.pdf}
\caption{The scalar-fiber Hamiltonian bracket closes in the resolved spectral regime.  Aliasing, a derivative without an exact product rule, and an incomplete detail constraint all produce large defects.}
\label{fig:constraint-negative-controls}
\end{figure}

These tests are stronger than substituting one favorable trajectory.  They show that the small residual is tied to the theorem's structural assumptions.

\section{Relation to the polynomial Lorentzian parent}

The first-order matter action \eqref{eq:cov-first-order} can be appended to the existing Einstein--Cartan parent before the $3+1$ split.  The normal tetrad multiplier then receives one additional scalar-fiber constraint, while tangential diffeomorphisms receive \eqref{eq:adm-dmatter} after regular reduction.  Level rotations commute with Lorentz transformations and spacetime diffeomorphisms.

At finite spectral resolution, the complete polynomial BV master identity transfers whenever every monomial generated by the action and antibracket lies below Nyquist.  A conservative headroom rule can be obtained by counting the eight-field potential monomial together with antifield variations.  This rule is sufficient rather than sharp.  Strict finite-range locality, a quantum BV measure, and uniform continuum estimates remain separate tasks.

\section{What has and has not been closed}

\begin{longtable}{@{}P{0.45\textwidth}P{0.19\textwidth}P{0.26\textwidth}@{}}
\toprule
Statement & Status & Boundary\\
\midrule
finite-depth covariant scalar-fiber action & closed & classical EFT\\
total stress conservation & closed & on shell\\
average-detail exchange law & closed & exact at finite depth\\
matter $H$--$H$ bracket and total ADM algebra & closed & regular classical chart with gravity included\\
resolved matter Hamiltonian bracket & closed conditional & alias-free spectral headroom\\
polynomial first-order parent & closed & nondegenerate tetrad chart\\
strict finite-range nonlinear realization & open & corner-to-continuum construction required\\
quantum master equation and anomaly & open & regulator and measure required\\
self-consistent production state & open & addressed diagnostically in \Cref{ch:backreaction}\\
\bottomrule
\end{longtable}
